\section{Worked BW500 wire exchange}
\label{app:bw500-wire}

This radio-free example follows a Linkwise caller \texttt{W9SSJ} sending
200 uncompressed bytes to \texttt{W1A}. Host byte $i$ has value $i$ for
$0\le i<200$, so the payload is the inclusive hexadecimal sequence
\texttt{00} through \texttt{c7}. It uses the measured frame layout exercised
by \path{modem::data::tests::reassembles_the_documented_200_byte_delivery_without_padding}.
The bytes below are reproducible codec vectors. The exchange diagram
combines those vectors with the implemented virtual-host session sequence.

\subsection{Establish the session}

The caller sends a request keyed to \texttt{W1A}. The responder returns a
response keyed to the same callsign. The caller then sends its identity
in a level-4 setup frame. Table~\ref{tab:worked-bw500-setup} gives every
byte before whitening and channel coding. The receiver checks the frame
CRC, decodes the identity, and sends a connected ACK keyed to \texttt{W9SSJ}.
The caller sends a confirm keyed to \texttt{W1A}, waits for the responder's
ACK, and obtains the transmit turn through the poll and turn-request exchange.

\begin{table}[htbp]
\centering\small
\begin{tabular}{@{}rll@{}}
\toprule
Offset & Hex bytes & Meaning \\
\midrule
0--4 & \texttt{5e 44 d3 28 00} & Packed six-bit callsign symbols and trailing zero bits \\
5--6 & \texttt{00 00} & No SSID and a reserved zero byte \\
7 & \texttt{80} & Base capability \\
8--9 & \texttt{14 f6} & Marker and high byte of the callsign CRC \\
10--41 & \texttt{00} repeated 32 times & Padding emitted by the setup encoder \\
42--43 & \texttt{04 82} & Requested level and setup control \\
44--45 & \texttt{1d ee} & Big-endian frame CRC over offsets 0--43 \\
\bottomrule
\end{tabular}
\caption{BW500 level-4 setup frame for \texttt{W9SSJ}. Offsets are decimal and zero-based.}
\label{tab:worked-bw500-setup}
\end{table}

The six-bit symbol values are $23,35,19,19,10,0$. The last value fills the
unused sixth character slot. CRC-16/GENIBUS of the uppercase ASCII callsign
is \texttt{f677}, so the setup and closing DATA trailer both use \texttt{f6}.
The frame CRC is separate and covers the body including that trailer.

\subsection{Build and decode all five DATA frames}

The level-4 frames each carry 43 payload bytes at offsets 0--42, one control
byte at offset 43, and a two-byte CRC at offsets 44--45. The fifth frame
uses level 3 and carries the remaining 28 bytes. The inclusive hex ranges
in Table~\ref{tab:worked-bw500-data} denote consecutive byte values, so each
row specifies the complete frame. The sender's \texttt{BUFFER} count is
shown after the corresponding acknowledgment.

\begin{table}[htbp]
\centering\small
\begin{tabular}{@{}rrlrrlr@{}}
\toprule
Frame & Level & Payload values & Bytes & Control & CRC & \texttt{BUFFER} \\
\midrule
1 & 4 & \texttt{00--2a} & 43 & \texttt{95} & \texttt{cf 4d} & 157 \\
2 & 4 & \texttt{2b--55} & 43 & \texttt{91} & \texttt{fc 74} & 114 \\
3 & 4 & \texttt{56--80} & 43 & \texttt{8d} & \texttt{e0 ea} & 71 \\
4 & 4 & \texttt{81--ab} & 43 & \texttt{81} & \texttt{17 49} & 28 \\
5 & 3 & \texttt{ac--c7} & 28 & \texttt{82} & \texttt{f0 fe} & 0 \\
\bottomrule
\end{tabular}
\caption{The 200-byte BW500 DATA delivery. The frame-5 trailer is shown below.}
\label{tab:worked-bw500-data}
\end{table}

Frame 5 contains the following 37 bytes, including its CRC.
\begin{verbatim}
00  ac ad ae af b0 b1 b2 b3 b4 b5 b6 b7 b8 b9 ba bb
10  bc bd be bf c0 c1 c2 c3 c4 c5 c6 c7 14 f6 00 00
20  00 04 82 f0 fe
\end{verbatim}
The left column gives hexadecimal offsets. In decimal, offsets 28 and 29
hold the marker \texttt{14} and caller key \texttt{f6}, offsets 30--32 hold
zero fill, and offset 33 advertises level 4 for the next receive over.
Offset 34 holds the closing control \texttt{82}. The CRC at offsets 35--36
covers offsets 0--34. The advertised level 4 is a field in the level-3 frame.

\landscapediagram{\begin{tikzpicture}[x=0.44cm,y=0.75cm,font=\small,
  field/.style={draw=ovnavy,thin},
  covered/.style={ovblue,decorate,decoration={brace,mirror,amplitude=4pt}}]
  \node[anchor=west,font=\bfseries] at (0,11.5) {Decoded BW500 frame bytes};
  \draw[ovgray] (0,10.8) -- (46,10.8);
  \foreach \offset in {0,10,20,30,40,46}{
    \draw[ovgray] (\offset,10.8) -- ++(0,0.13)
      node[above,font=\scriptsize] {\offset};
    \draw[ovgray!23] (\offset,10.8) -- (\offset,4.9);
  }
  \node[anchor=east,font=\scriptsize] at (46,11.5) {Zero-based byte boundaries};

  \node[anchor=west,font=\bfseries\small] at (0,10.25) {Level-4 setup \quad 46 bytes};
  \foreach \from/\to/\shade/\name in {
    0/10/ovteal!18/{Identity header},
    10/42/ovpale/{Zero fill},
    42/44/ovgold!65/{},
    44/46/ovblue!25/{}} {
    \draw[field,fill=\shade] (\from,9) rectangle (\to,9.8);
    \node at ({(\from+\to)/2},9.4) {\name};
  }
  \draw[covered] (0,8.9) -- (44,8.9)
    node[midway,below=5pt,font=\scriptsize] {CRC covers bytes 0--43};
  \node[anchor=east,font=\scriptsize] at (46,8.25)
    {\texttt{04 82} at 42--43 \quad CRC \texttt{1d ee} at 44--45};

  \node[anchor=west,font=\bfseries\small] at (0,7.8) {Full level-4 DATA \quad 46 bytes};
  \draw[field,fill=ovteal!18] (0,6.55) rectangle (43,7.35)
    node[midway] {43 payload bytes};
  \draw[field,fill=ovgold!65] (43,6.55) rectangle (44,7.35);
  \draw[field,fill=ovblue!25] (44,6.55) rectangle (46,7.35);
  \draw[covered] (0,6.45) -- (44,6.45)
    node[midway,below=5pt,font=\scriptsize] {CRC covers bytes 0--43};
  \node[anchor=east,font=\scriptsize] at (46,5.8)
    {Frame 1 control \texttt{95} at 43 \quad CRC \texttt{cf 4d} at 44--45};

  \node[anchor=west,font=\bfseries\small] at (0,5.35) {Closing level-3 DATA \quad 37 bytes};
  \draw[field,fill=ovteal!18] (0,4.1) rectangle (28,4.9)
    node[midway] {28 payload bytes \quad \texttt{ac--c7}};
  \foreach \from/\to/\shade in {28/30/ovgold!65,30/33/ovpale,33/35/ovgold!65,35/37/ovblue!25}
    \draw[field,fill=\shade] (\from,4.1) rectangle (\to,4.9);
  \draw[covered] (0,4.0) -- (35,4.0)
    node[midway,below=5pt,font=\scriptsize] {CRC covers bytes 0--34};
  \node[anchor=west,font=\bfseries\small] at (4,2.8) {Closing trailer and CRC expanded};
  \foreach \offset/\value/\shade [count=\i from 0] in {
    28/14/ovgold!65,29/f6/ovgold!65,30/00/ovpale,31/00/ovpale,
    32/00/ovpale,33/04/ovgold!65,34/82/ovgold!65,35/f0/ovblue!25,36/fe/ovblue!25} {
    \draw[field,fill=\shade] ({4+\i*4.2},1.15) rectangle ++(4.2,0.85);
    \node[font=\ttfamily] at ({6.1+\i*4.2},1.58) {\value};
    \node[font=\scriptsize] at ({6.1+\i*4.2},2.22) {\offset};
  }
  \foreach \x/\name in {6.1/Marker,10.3/Caller key,18.7/Zero fill,27.1/Next level,31.3/Control,37.6/CRC}
    \node[font=\scriptsize] at (\x,0.76) {\name};
\end{tikzpicture}
}{Byte offsets and CRC coverage for the worked setup, full DATA, and closing DATA frames. The three rows share a byte scale. The inset expands the closing trailer and CRC.}{fig:worked-bw500-layout}

The transmitter serializes each frame most-significant bit first. At
level 4, 46 bytes give 368 information bits and the whitening, turbo and
interleaving chain yields 748 bits for the differential cells. At level 3,
37 bytes give 296 bits. One zero pad bit outside the byte CRC brings the
information block to 297 bits, yielding 606 coded bits mapped to 202 tone-index
columns. The receiver reverses the matching waveform and coding pipeline,
validates the frame CRC, and interprets the DATA control and trailer.
It buffers the first 172 bytes and releases all 200 to the host only when
the closing frame validates. The marker, key, padding, control, and CRC
are excluded from the host stream.

For an independent byte check, the following Python reproduces all six
frames, including the setup frame. \texttt{crc\_hqx} computes the
unreflected polynomial update and the final XOR supplies the complement.
\begin{verbatim}
from binascii import crc_hqx

def crc(body):
    return (crc_hqx(body, 0xffff) ^ 0xffff).to_bytes(2, "big")

payload = bytes(range(200))
key = crc(b"W9SSJ")[:1]
setup = bytes.fromhex("5e 44 d3 28 00 00 00 80 14")
setup += key + bytes(32) + bytes.fromhex("04 82")
bodies = [setup]
for index, control in enumerate((0x95, 0x91, 0x8d, 0x81)):
    start = 43 * index
    bodies.append(payload[start:start + 43] + bytes([control]))
bodies.append(payload[172:] + b"\x14" + key + bytes(3) + b"\x04\x82")
for body in bodies:
    print((body + crc(body)).hex(" "))
\end{verbatim}

\subsection{Trace frame 1 through the physical layer}
\label{sec:worked-bw500-phy}

Take the first DATA frame from Table~\ref{tab:worked-bw500-data}.
Its 46 bytes are \texttt{00} through \texttt{2a}, followed by
\texttt{95 cf 4d}. This trace uses the checked-in
\path{reference/vara/fixtures/level4/bw500-level4.json} profile and
48 kHz mono PCM from \path{modem::bw500_waveform::Waveform::synthesize}.
The receiver sees the complete synthetic burst with nominal carrier and
sample rates, unit gain, zero added noise, and no channel filtering.
It searches for timing within that candidate. The example characterizes
the transforms and framing on this input. The profile SHA-256 is
\path{7fe1eed44ca50f20cb7f64e1e70b9e21d8fe9c348a920609b9888bfdde82bb48}.

\paragraph{Serialization, whitening, and turbo coding}
Serialize the bytes most-significant bit first. The first 16 input bits
are \texttt{00000000 00000001}. The whitening generator starts at 29,318.
Its first 16 output bits are \texttt{11101010 11111110}, giving
\texttt{11101010 11111111} after XOR. The full 368 whitened information
bits, packed most-significant bit first, are
\par\noindent\begin{minipage}{\linewidth}
\begin{verbatim}
ea ff 1a 2d 5a 33 09 a8 83 9c 53 d3 f8 0d 52 32
fe ee 28 a6 b0 49 b4 83 68 f4 fc a5 2c e9 d2 d9
8b e4 c8 c2 a8 a4 8a 15 c5 62 50 70 55 4c
\end{verbatim}
\end{minipage}

With zero-based indexing, the second constituent starts with whitened
bits at indices 271, 94, 29, 148, 87, 280, 290, and 316 in that order. The encoder
pairs each systematic bit with the first constituent's parity at even
indices and the second's parity at odd indices. Three termination steps
per constituent contribute the final 12 bits. The complete 748-bit
vector before channel interleaving is
\begin{verbatim}
fd cc bf ea 16 99 0c e3 37 9c 1b 1e 14 93 88 95
d5 4b d2 b5 63 1a b3 5b fb 91 44 a3 23 5c 4a 48
ef fc a9 b9 19 c0 98 38 8b 15 21 c3 8f 24 d4 5b
2d d1 ff 64 ba b1 d9 62 19 e4 ec 93 e3 4d a7 87
c1 cb fd 30 b4 c0 a0 0d 88 c5 9c 61 d4 cc 02 73
f5 36 7d 4c 33 01 7a 10 37 33 24 e1 c1 30
\end{verbatim}
In each 748-bit hex dump, the low nibble of the last byte is zero
display padding. The termination bits are \texttt{110000010011}.

\paragraph{Channel order and symbol placement}
The channel permutation gathers coded bits at indices
33, 165, 300, 126, 467, 684, 192, and 25 for the first eight transmitted positions.
The resulting channel-order vector is
\begin{verbatim}
03 0b 87 87 f8 33 2b 92 a7 5d 99 23 01 58 b6 8c
5b dc 7e 6e d9 38 dc 42 41 52 e4 3e 33 d9 63 a9
1f 2d 03 0b 52 bc 8a ec 44 fa 0d 69 82 f1 48 2a
6a c3 e8 f8 56 b2 b0 5b e9 4e 96 2d 55 1b 6b 2e
06 bd 01 1c 78 6c 5f b9 c5 aa 7a b6 3f a1 4b 24
ef 1e 42 e6 ac f6 f8 28 14 54 c0 d4 d7 10
\end{verbatim}

Figure~\ref{fig:dsp-level4-symbol-timing} gives the transmit centers
$n_{b,j}=4608+512j+256b$. Column zero initializes the differential
polarity in each band. Bits $q_0$ through $q_{373}$ fill the lower band
and $q_{374}$ through $q_{747}$ fill the upper band. In each band the
mapper skips pilot columns. The first six lower-band DATA bits leave the
polarity unchanged. Bits $q_6$ and $q_7$ flip it on successive columns.
For example, $q_6$ lands at lower-band column 7. Start at sample 4608
and add seven 512-sample hops to get sample 8192. The upper-band center
at the same column comes 256 samples later.
In the upper band, $q_{379}$ occupies column 6, column 7 is a pilot,
and $q_{380}$ occupies column 8. The pilot holds polarity and consumes
no coded bit.

\paragraph{Timing estimate, soft metrics, and recovered frame}
For the 205,825-sample synthesized burst, the timing search selects
$\widehat n_0=4524$ with a mean normalized pilot score of approximately
0.99989, above the 0.65 threshold. The nominal transmit center is 4608.
The estimated epoch maximizes pilot agreement for this waveform and
search window. Its 84-sample difference is specific to this trace.

At each estimated center $\widehat n_{b,j}=\widehat n_0+512j+256b$,
the receiver mixes to complex baseband, sums against the 513 pulse taps,
and removes the measured reference phase. Call the resulting complex
sample $z_{b,j}$. For each DATA column it computes
$d_{b,j}=\Re\{z_{b,j}\overline{z_{b,j-1}}\}$ and emits
$m_i=-6d_{b,j}/\overline{|d|}$. In this trace
$\overline{|d|}\simeq408.10681$ across the 748 DATA cells.
Here $\overline{|d|}$ is the average size of the dot products with their
signs removed. Dividing by that average puts the cells on a common
scale. The factor of six sets their average absolute score to six.
For $q_6$, the dot product is about $-406.93$. Negating and scaling
it gives $+5.98269$, evidence for a polarity flip and hence bit 1.
Table~\ref{tab:worked-bw500-metrics} shows selected positions and their
rounded metrics. The preceding column supplies the reference for each
differential decision, including the upper-band pilot immediately before
column 8.

\begin{table}[htbp]
\centering\small
\begin{tabular}{@{}rrrrrrrr@{}}
\toprule
$i$ & Coded index & $q_i$ & Band & Column & TX center & RX center & $m_i$ \\
\midrule
0   & 33  & 0 & 0 & 1   & 5120   & 5036   & $-5.88364$ \\
6   & 192 & 1 & 0 & 7   & 8192   & 8108   & $+5.98269$ \\
7   & 25  & 1 & 0 & 8   & 8704   & 8620   & $+5.93823$ \\
373 & 440 & 0 & 0 & 392 & 205312 & 205228 & $-5.98740$ \\
374 & 733 & 0 & 1 & 1   & 5376   & 5292   & $-6.15735$ \\
379 & 293 & 0 & 1 & 6   & 7936   & 7852   & $-6.12393$ \\
380 & 457 & 1 & 1 & 8   & 8960   & 8876   & $+5.95272$ \\
\bottomrule
\end{tabular}
\caption{Selected channel-order soft metrics for frame 1. Centers are absolute sample indices within the synthetic burst.}
\label{tab:worked-bw500-metrics}
\end{table}

For example, scatter $m_6$ to coded index 192. Index 192 is the
systematic position for information bit 96, so its positive sign yields
whitened bit 1 at that position. The receiver scatters the remaining
metrics the same way and takes the signs of coded positions
$0,2,\ldots,734$. The systematic decisions reproduce all 368 whitened
bits. Applying the whitening sequence again recovers
\texttt{00} through \texttt{2a} followed by \texttt{95 cf 4d}.
CRC-16/GENIBUS over the first 44 bytes yields \texttt{cf4d}.
The clean trace therefore exits on the systematic CRC check.
If that check fails, the iterative max-log-MAP path in
\path{modem::coding} uses the parity metrics before trying the CRC again.
The modem contributes 43 bytes to the pending delivery from this frame.

\subsection{Confirm delivery and recover a withheld ACK}

An accepted intermediate reply advances the sender to the next DATA frame.
The closing ACK confirms the delivery at the modem boundary and retires
the last 28 bytes. If the closing ACK is withheld, the caller still reports
28 buffered bytes. It sends a short final query keyed to the called station
and can follow with the long final query. A Linkwise receiver in the closing
ACK window answers with the same caller-keyed ACK without emitting the
payload a second time. After a valid ACK the caller reports \texttt{BUFFER 0}.
A subsequent host disconnect sends the session disconnect control.
This example ends at modem DATA delivery. Winlink mailbox storage is a
separate test outcome.

\landscapeplantuml{build/diagrams/Linkwise_BW500_Worked_Exchange}{BW500 setup, transmit-turn acquisition, DATA acknowledgments, and closing-ACK recovery for the worked delivery. The recovery branch shows the Linkwise virtual-host behavior exercised by the socket fixture.}{fig:worked-bw500-exchange}

\subsection{Other implemented frame layouts}

Table~\ref{tab:worked-mode-widths} gives the decoded byte geometry. The
control offset is also the final byte covered by the frame CRC.
Full frames carry payload in every byte before control. A keyed closing
frame uses a marker, caller key, optional zero fill, and rate byte before
control. The decoder also accepts measured BW500 closing variants.
The keyed forms use \texttt{42} or \texttt{82}, with \texttt{02} accepted
at level 4. A 32-byte level-3 or 41-byte level-4 close can use
\texttt{14} followed directly by the rate and \texttt{82}. An unpadded
\texttt{83} close carries 33 bytes at level 3 or 42 at level 4, followed
by the rate and control. The encoder emits the keyed form by default.
It also emits the two delimiter-only forms and the 42-byte unpadded
level-4 form.

\begin{table}[htbp]
\centering\small
\begin{tabular}{@{}lrrrrl@{}}
\toprule
Mode & Level & Full payload & Control offset & Frame bytes & CRC offsets \\
\midrule
BW500 & 1 & 9 & 9 & 12 & 10--11 \\
BW500 & 2 & 22 & 22 & 25 & 23--24 \\
BW500 & 3 & 34 & 34 & 37 & 35--36 \\
BW500 & 4 & 43 & 43 & 46 & 44--45 \\
BW2300 / BW2750 & 2 & 22 & 22 & 25 & 23--24 \\
BW2300 / BW2750 & 3 & 47 & 47 & 50 & 48--49 \\
BW2300 / BW2750 & 4 & 89 & 89 & 92 & 90--91 \\
\bottomrule
\end{tabular}
\caption{Decoded DATA byte geometry before mode-specific coding and modulation.}
\label{tab:worked-mode-widths}
\end{table}

Wide level-4 setup keeps the same identity header at offsets 0--9. Its
encoder fills offsets 10--87 with zeros, places the level and control at
88--89, and appends CRC at 90--91. BW2300 and BW2750 share the byte
geometry in the table and use different carrier allocations. The BW500
example's control sequence and waveform parameters apply to its tested
profile. The wide receive and recovery results are tracked separately in
Table~\ref{tab:modem-evidence-gates}.
